\documentclass[conference]{IEEEtran}
\IEEEoverridecommandlockouts

\usepackage{cite}
\usepackage{amsmath,amssymb,amsfonts}
\usepackage{algorithmic}
\usepackage{graphicx}
\usepackage{textcomp}
\usepackage{xcolor}
\usepackage{booktabs}
\usepackage{url}
\usepackage{verbatim}
% Resolves whether the figure sits beside paper.tex (Overleaf upload of paper/
% alone) or in the repo's figures/ directory.
\graphicspath{{./}{../figures/}}

\def\BibTeX{{\rm B\kern-.05em{\sc i\kern-.025em b}\kern-.08em T\kern-.1667em\lower.7ex\hbox{E}\kern-.125emX}}

\newtheorem{theorem}{Theorem}
\newtheorem{lemma}{Lemma}
\newtheorem{proposition}{Proposition}
\newtheorem{definition}{Definition}

\begin{document}

\title{Anchored Counterfactual Attribution of Model\\Degradation to ETL Pipeline Operators}

\author{
\IEEEauthorblockN{Laiba Mazhar}
\IEEEauthorblockA{\textit{Data Scientist}\\
Lahore, Pakistan\\
lm4442172@gmail.com}
\and
\IEEEauthorblockN{Maryam Sarfraz}
\IEEEauthorblockA{\textit{Computer Scientist}\\
Lahore, Pakistan\\
sarfrazmaryam123@gmail.com}
}

\maketitle

\begin{abstract}
Production machine learning models are fed by ETL pipelines of tens to hundreds
of operators. When model quality drops between two scheduled runs with no code
deployed, engineers must determine which operator is responsible. Current
practice walks the pipeline in topological order, swapping in the new run's data
one stage at a time. We show this is the marginal-contribution vector of a single
permutation and is therefore order-dependent: on our benchmark the \emph{sign} of
attributed blame reverses between two valid topological orders of the same DAG.
We formulate the problem as a cooperative game. Holding the pipeline fixed, each
operator executes in either of two states---as in a known-good \emph{reference}
run or as in the degraded \emph{incident} run---and a \emph{hybrid replay}
executes the DAG under an arbitrary state assignment. The Shapley value of the
resulting game decomposes the observed degradation exactly. Our central finding
is negative and then constructive: the Shapley value is correct as accounting but
poor as triage, ranking the true culprit first in only 60\% of measurable
incidents, and failing precisely when the injected fault does less damage than
co-occurring benign drift. Leave-one-out ablation ranks well on synthetic data
but its attributions deviate from the observed degradation by up to $0.246$~AUC.
We introduce the \emph{incident-anchored Shapley value}, which holds a designated
set of exogenous operators at incident state and takes the Shapley value of the
resulting sub-game. It recovers standard Shapley and leave-one-out as boundary
cases, attains 100\% top-1 accuracy in every severity regime, and retains exact
decomposition to within $3\times10^{-17}$. On real NYC taxi partitions the
separation sharpens: comparing a Monday against a Saturday, leave-one-out
collapses to $0.50$ while the anchored value holds at $1.00$. Output-hash pruning
and prefix memoisation reduce $2^{33}$ hybrid replays to three model fits on a
33-operator DAG.
\end{abstract}

\begin{IEEEkeywords}
data pipelines, model degradation, root cause analysis, Shapley value, data-centric AI, ETL
\end{IEEEkeywords}

\section{Introduction}

A model scores $0.87$ AUC on Monday. On Tuesday the same Airflow DAG runs,
nothing was deployed, and it scores $0.81$. Which task caused it?

This is among the most common failure modes in production machine learning, and
existing tooling addresses it poorly. Monitoring systems~\cite{deequ} report
\emph{that} a distribution moved; in a deep DAG a single upstream change makes
every downstream node move, so a monitor fires everywhere and ranks nothing.
Lineage systems~\cite{openlineage} report what is upstream of what, narrowing the
suspects to ``everything upstream.'' Pipeline debuggers---\textsc{mlinspect}
\cite{mlinspect,mlinspectj} is the strongest---instrument a pipeline into a
dataflow DAG and propagate annotations to detect data-distribution bugs
\emph{within one execution}; they do not compare executions across time and do
not assign quantitative blame.

So the work is done by hand, and the manual procedure has a defect nobody has
named: it depends on the order the engineer happens to walk the DAG in.

\subsection{Contributions}

\begin{enumerate}
\item \textbf{Formulation} (Sec.~\ref{sec:game}). Temporal hybrid replay of a
fixed pipeline over two snapshots as a cooperative game whose Shapley value
decomposes observed degradation exactly. Verified numerically to
$6\times10^{-17}$.
\item \textbf{A negative result} (Sec.~\ref{sec:negative}). The Shapley value is
the wrong estimator for triage, failing when the fault is smaller than
co-occurring benign drift.
\item \textbf{The anchored Shapley value} (Sec.~\ref{sec:anchored}). A
one-parameter family spanning standard Shapley and leave-one-out, attaining
100\% top-1 accuracy while preserving exact decomposition.
\item \textbf{Tractability} (Sec.~\ref{sec:computing}). Output-hash pruning and
prefix memoisation, with the collapse regime identified.
\item \textbf{A fault-injection benchmark} over synthetic and real NYC taxi
pipelines, with ground-truth culprits, monotone severity ladders, and
\emph{benign controls}---faults that violate data-quality constraints loudly
while causing no model damage.
\item \textbf{A real-data result} (Sec.~\ref{sec:real}) showing that
leave-one-out, indistinguishable from the anchored value on synthetic data,
fails under real weekday/weekend drift.
\end{enumerate}

We do not have a production deployment, we have not integrated with an
orchestrator, our real-data evaluation is one dataset and one pipeline topology,
and the anchored estimator's axiomatic characterisation is stated but not proven.
Section~\ref{sec:limits} is explicit about all of this.

\section{Problem Statement}

A pipeline is a DAG $G=(V,E)$ of operators, each computing
$d_i = f_i(\text{inputs}, \text{params}_i)$. Two executions are given: a
\emph{reference} run $R$ (last known-good) and an \emph{incident} run $I$
(degraded). Between them each operator may differ in its source partition, its
configuration, or its code. We encode all three uniformly as an operator
\emph{state} $\sigma_i \in \{R, I\}$.

A downstream model is trained on the sink dataset and evaluated on a fixed probe
set. Writing $u$ for that utility, the observed degradation is
$\Delta = u(V) - u(\emptyset)$.

\textbf{The attribution problem.} Distribute $\Delta$ across operators such that
the assignment identifies which operators to repair.

\section{Hybrid Replay and the Degradation Game}\label{sec:game}

\begin{definition}[Hybrid replay]
For $S \subseteq V$, $\Pi(S)$ executes the DAG with every operator in $S$ in
incident state and every operator in $V \setminus S$ in reference state:
\begin{equation}
d_i(S) = f_i\!\left(\{d_j(S) : (o_j,o_i)\in E\},\ \text{params}_i^{\sigma_i(S)}\right)
\end{equation}
where $\sigma_i(S) = I$ if $o_i \in S$ and $R$ otherwise.
\end{definition}

$\Pi(\emptyset)$ reproduces the good run, $\Pi(V)$ the bad one, and the $2^n-2$
hybrids between them interpolate.

\begin{definition}[Degradation game]
With $u(S)$ the model utility on $\Pi(S)$'s output, $v(S) = u(S) - u(\emptyset)$.
\end{definition}

Then $v(\emptyset)=0$ and $v(V)=\Delta$. The probe set is held fixed across every
coalition; otherwise $u(S)$ compares different things and every property below is
vacuous. The Shapley value
\begin{equation}
\phi_i = \sum_{S \subseteq V\setminus\{i\}} \frac{|S|!\,(n-|S|-1)!}{n!}
         \left[v(S\cup\{i\}) - v(S)\right]
\label{eq:shapley}
\end{equation}
satisfies \emph{efficiency}: $\sum_i \phi_i = \Delta$. The degradation is
\emph{decomposed}, not merely scored. No baseline we are aware of has this
property.

\subsection{Current Practice is a Single Permutation}

Walking the DAG in topological order and swapping stages in one at a time
produces, for order $\pi$, the vector $m_i^\pi = v(S_\pi(i)\cup\{i\}) - v(S_\pi(i))$
where $S_\pi(i)$ precedes $i$. This is the marginal-contribution vector of
\emph{one} permutation; \eqref{eq:shapley} averages over all $n!$.

\begin{proposition}
Single-permutation stagewise diagnosis violates the symmetry axiom and is
order-dependent.
\end{proposition}

The empirical form is stark. Enumerating \emph{every} valid topological order of
the seven-operator pipeline over four fault scenarios,
\textbf{8 of 28 operator/scenario pairs receive blame of both signs} depending
only on the order chosen (Table~\ref{tab:order}). For \texttt{load\_customers}
the assigned blame ranges from $-0.0556$ to $+0.0405$: a spread of $0.096$~AUC on
incidents whose total magnitude is $0.03$--$0.11$~AUC. Two engineers running the
identical manual procedure on the identical incident reach opposite conclusions
about whether an operator helped or harmed.

\begin{table}[t]
\caption{Blame under all topological orders (7-operator pipeline)}
\label{tab:order}
\centering
\begin{tabular}{lccc}
\toprule
Scenario & Orders & Stagewise top-blamed & Anchored \\
\midrule
unit change      & 6 & \texttt{agg\_transactions} (6/6) & correct \\
predicate flip   & 6 & \texttt{filter\_active} (6/6)    & correct \\
schema drift     & 6 & \emph{a source} (6/6, wrong)     & correct \\
unit + predicate & 6 & \texttt{agg\_transactions} (5/6) & correct \\
\bottomrule
\end{tabular}
\end{table}

Worse than noisy: on the schema-drift scenario stagewise diagnosis blames an
innocent \emph{source} operator in all six orders, never the culprit. The
anchored value is order-independent by construction and is correct in all four
(Fig.~\ref{fig:order}).

\begin{figure}[t]
\centering
\includegraphics[width=\columnwidth]{order_dependence.pdf}
\caption{Range of blame stagewise diagnosis assigns to each operator across all
six topological orders of the unit-change scenario. Bars crossing zero change
sign with the order alone. The anchored value is a single point regardless of
order.}
\label{fig:order}
\end{figure}

\section{The Incident-Anchored Shapley Value}\label{sec:anchored}

Section~\ref{sec:negative} shows the Shapley value ranks culprits worse than
naive leave-one-out. The reason is a conditioning difference. Leave-one-out
evaluates $v(V)-v(V\setminus\{i\})$, the marginal effect \emph{at the actual
incident configuration}, where benign drift is held fixed and cancels. Shapley
averages over all $2^n$ configurations, including many where the benign drift has
not happened, so drift receives a real, nonzero share. Both are right about
different questions, and the space between them is empty in the literature.

\begin{definition}[Incident-anchored Shapley value]
For an anchor set $B \subseteq V$, let $v_B(T) = v(T \cup B) - v(B)$ be the
sub-game on $V\setminus B$. The anchored value $\phi^B$ is the Shapley value of
$v_B$, with $\phi_i^B = 0$ for $i \in B$ and, writing $m = |V \setminus B|$,
\begin{equation}
\phi_i^B \!=\!\! \sum_{S \subseteq V\setminus(B\cup\{i\})}\!\!\!
  \frac{|S|!\,(m-|S|-1)!}{m!}
  \left[v(S\cup B\cup\{i\}) - v(S\cup B)\right]
\end{equation}
\end{definition}

\begin{theorem}[Boundary cases]
$B=\emptyset$ gives the standard Shapley value. $B = V\setminus\{i\}$ gives
$\phi_i^B = v(V) - v(V\setminus\{i\})$, the leave-one-out value of $i$.
\end{theorem}

Both are verified as assertions to machine precision before any result is
reported.

\begin{theorem}[Anchored efficiency]
$\sum_{i \notin B} \phi_i^B = v(V) - v(B)$.
\end{theorem}

The anchored values decompose exactly the part of the degradation not already
explained by the anchor. Measured maximum deviation across all 29 synthetic
configurations: $2.78\times10^{-17}$.

\subsection{Choosing the Anchor}

The principled default anchors the \emph{exogenous} operators---sources, with no
parents. Their state differs because the world moved, not because the pipeline
did. One cannot repair the fact that Tuesday's data is not Monday's, so charging
blame to it is accurate accounting and useless triage. This is a modelling
decision and we present it as one.

\section{Computing It}\label{sec:computing}

Each $v(S)$ costs a pipeline replay plus a model fit; $2^n$ of them is hopeless.
Three mechanisms make it practical.

\textbf{Ancestral sufficiency.} $u(S)$ depends only on the sink, so operators
outside $\mathrm{Anc}^*(\text{sink})$ are null players. In real Airflow DAGs most
tasks are sensors, notifications and cleanup; this removes them for free.

\textbf{Output-stability pruning.} If operator $i$'s output is byte-identical
under both states given the same inputs, all downstream computation is identical
and $v(S\cup\{i\}) = v(S)$. This is detectable by \emph{hashing}, with no model
fit---one hash replaces one replay-and-retrain. The operators surviving this test
form the \emph{active frontier} $F$.

\textbf{Prefix memoisation.} Hybrid replays share large prefixes. We cache
operator outputs on $(\text{operator}, \text{state}, \text{hash(inputs)})$ and
utilities on the sink hash, so distinct coalitions producing an identical sink
share a single model fit.

Beyond $n \approx 14$ we use permutation sampling over the free players with the
anchor always included.

\section{Evaluation}

\subsection{Setup}

\textbf{Synthetic pipeline.} A seven-operator churn feature build: two sources, a
cleaning step, an aggregation, a join, a filter, an encoder; 4000 customers per
day. $n=7$ keeps the 128-coalition lattice brute-forceable so Monte-Carlo
estimates can be checked against exact ground truth.

\textbf{Benign drift.} Reference and incident runs read \emph{different days} of
source data from the same distribution, so every scenario layers a real fault on
natural churn and every downstream node's output changes whether or not it is to
blame. This is what makes per-node monitoring fail, and it is what breaks the
plain Shapley value.

\textbf{Real pipeline.} Fig.~\ref{fig:pipeline} shows the nine-operator NYC taxi
pipeline used in Sec.~\ref{sec:real}, with the two operators carrying injected
faults highlighted. CULPA is not told which they are.

\begin{figure*}[t]
\centering
\includegraphics[width=\textwidth]{pipeline_taxi.pdf}
\caption{The nine-operator NYC yellow taxi pipeline. Sources (left) are
exogenous---their state differs because the world moved between the two runs, not
because the pipeline did---so they form the default anchor set. The two
highlighted operators carry the injected faults of Fig.~\ref{fig:attribution}.}
\label{fig:pipeline}
\end{figure*}

\textbf{Faults.} Five types with severity ladders (29 configurations): unit
change, schema drift, filter-predicate flip, MNAR null spike, biased join
fan-out. Two \emph{benign controls}---uniform join fan-out (35\% duplicate keys)
and MCAR nullness (55\% of rows dropped)---violate data-quality constraints
loudly while causing no model damage.

Fig.~\ref{fig:severity} plots what each fault actually costs as severity rises,
and it is why the evaluation reports exclusions. Several configurations sit below
any measurable effect, and one ladder is not even monotone: for schema drift
\emph{which} column is dropped matters more than how many, so losing
\texttt{tenure} costs an order of magnitude more than losing
\texttt{total\_amount}. An injected fault is a hypothesis about damage and has to
be checked before attribution is scored against it.

\begin{figure}[t]
\centering
\includegraphics[width=\columnwidth]{severity_ladder.pdf}
\caption{Damage each fault causes in isolation, by severity level. Below the
dashed threshold there is no degradation to attribute and no culprit to find.
Schema drift is non-monotone: the identity of the dropped column dominates its
count.}
\label{fig:severity}
\end{figure}

\textbf{Baselines.} Per-node distributional drift, single-permutation stagewise
swap (current practice), and leave-one-out.

\subsection{The Negative Result}\label{sec:negative}

On the initial gate experiment (10 scenarios), leave-one-out beat the Shapley
value, as shown in Table~\ref{tab:gate}.

\begin{table}[t]
\caption{Initial gate experiment, 8 harmful scenarios}
\label{tab:gate}
\centering
\begin{tabular}{lccc}
\toprule
Method & p@1 & MRR & Blame mass \\
\midrule
Leave-one-out   & \textbf{1.00} & \textbf{1.00} & \textbf{0.87} \\
Per-node drift  & 0.88 & 0.94 & 0.51 \\
Shapley (exact) & 0.75 & 0.88 & 0.62 \\
Stagewise       & 0.38 & 0.69 & 0.41 \\
\bottomrule
\end{tabular}
\end{table}

Shapley lost exactly the two scenarios where the fault did \emph{less} damage
than benign drift: biased join fan-out ($\Delta = -0.0016$) and MNAR null spike
($\Delta = -0.0028$), against a per-source benign drift magnitude of $0.0556$. In
both it ranked a source above the culprit---which, as accounting, is correct.
This motivated the anchored value.

\subsection{Accuracy Against the Fault-to-Drift Ratio}

Let $r = |v(\{\text{culprit}\})| / \max_s |v(\{s\})|$ over sources $s$.
Figure~\ref{fig:main} plots precision@1 against $r$ for both pipelines.

\begin{figure*}[t]
\centering
\includegraphics[width=\textwidth]{main_result.pdf}
\caption{Precision@1 by fault-to-drift ratio $r$. The anchored value is flat at
$1.00$ in every regime. On the synthetic pipeline (a), plain Shapley degrades
below $r=1$ while leave-one-out does not. On real NYC taxi partitions (b),
leave-one-out collapses alongside Shapley and current practice, and the anchored
value is the only estimator that survives.}
\label{fig:main}
\end{figure*}

Across the 15 measurable synthetic configurations the anchored value attains
p@1 $=1.00$, against $0.60$ for plain Shapley, $0.87$ for per-node drift and
$0.33$ for stagewise. The predicted mechanism is confirmed: plain Shapley
degrades to $0.20$ when the fault is small relative to benign drift and recovers
to $1.00$ when it dominates, while the anchored value is flat because anchoring
removes benign drift by construction.

\subsection{Anchored Value versus Leave-One-Out}

On synthetic data the two rank identically---15/15 agreement on p@1. They differ
in what the numbers mean (Table~\ref{tab:efficiency}). Leave-one-out's
attributions miss the observed degradation by up to $0.246$~AUC, larger than most
of the incidents being diagnosed: it double-counts interacting faults, and in one
scenario reported that an operator \emph{helped} by $+0.04$ when its true
anchored contribution was $+0.005$.

\begin{table}[t]
\caption{Only the anchored value decomposes the degradation}
\label{tab:efficiency}
\centering
\begin{tabular}{lcc}
\toprule
 & Max deviation & Mean \\
\midrule
Anchored, $|\sum\phi - (v(V)-v(B))|$ & $\mathbf{2.78\times10^{-17}}$ & --- \\
Leave-one-out, $|\sum\phi - v(V)|$   & $2.46\times10^{-1}$ & $9.6\times10^{-2}$ \\
\bottomrule
\end{tabular}
\end{table}

\subsection{A Worked Incident}

Fig.~\ref{fig:attribution} shows the tool's output on a compound incident in the
real taxi pipeline of Fig.~\ref{fig:pipeline}: a filter predicate silently flips
and, independently, three feature columns vanish after an upstream rename. The
run loses $0.0379$~AUC. Anchoring the two sources removes the benign
weekday/weekend drift, leaving $0.0316$ to attribute, which the anchored value
splits $-0.0171$ to \texttt{filter\_valid} and $-0.0145$ to \texttt{encode} with
a residual of $2\times10^{-17}$. The remaining seven operators receive exactly
zero, having been eliminated by output-hash pruning without a single model fit.

That split is the operational point. A method that only ranks says
\texttt{filter\_valid} is worse; the decomposition says fixing it alone recovers
just over half the loss, which is what decides whether one repair is enough.

\begin{figure}[t]
\centering
\includegraphics[width=\columnwidth]{attribution.pdf}
\caption{Blame assigned on a compound incident in the taxi pipeline. Both
culprits are identified and the attributions sum to the degradation being
explained.}
\label{fig:attribution}
\end{figure}

\subsection{Benign Controls}

Table~\ref{tab:benign} contrasts the benign controls with a genuinely harmful
fault. Uniform fan-out and MCAR nullness fire every data-quality constraint a
monitor has---35\% duplicate keys, 55\% of rows dropped---and cost nothing. The
method assigns them near-zero blame.
\textbf{Not every data-quality violation is an incident}, and a constraint
checker cannot tell the two apart because it never looks at the model.

\begin{table}[t]
\caption{Constraint violation is not model damage}
\label{tab:benign}
\centering
\begin{tabular}{lccc}
\toprule
Injected fault & Constraint & $\Delta$ (AUC) & $\max|\phi|$ \\
\midrule
Uniform join fan-out, 35\% & fires & $+0.0001$ & 0.0078 \\
MCAR nullness, 55\% rows   & fires & $+0.0006$ & 0.0085 \\
Silent unit change         & silent & $-0.1149$ & 0.0996 \\
\bottomrule
\end{tabular}
\end{table}

\subsection{Real Data: NYC Yellow Taxi}\label{sec:real}

A nine-operator pipeline over NYC TLC yellow taxi records for January 2024: two
sources, a validity filter, time-feature derivation, a per-zone aggregate, two
zone joins, a stats join, and an encoder. Task: predict a generous tip
($\text{tip}/\text{fare} > 0.20$) for credit-card trips; 60{,}000 trips per
partition. \texttt{tip\_amount} and \texttt{total\_amount} are dropped before
modelling---the latter includes the tip and would leak the label outright.

Reference and incident runs read two different real days. Table~\ref{tab:real}
gives two regimes.

\begin{table}[t]
\caption{Real NYC taxi, precision@1 over valid test cases}
\label{tab:real}
\centering
\begin{tabular}{lcc}
\toprule
 & Mon 8 vs Tue 9 & Mon 8 vs Sat 13 \\
\midrule
Benign drift $\max|v(s)|$ & 0.0104 & 0.0063 \\
Valid test cases          & 6 / 19 & 8 / 19 \\
\midrule
Anchored Shapley          & \textbf{1.00} & \textbf{1.00} \\
Per-node drift            & 0.50 & 0.62 \\
Shapley                   & 1.00 & \textbf{0.50} \\
Leave-one-out             & 1.00 & \textbf{0.50} \\
Stagewise                 & 1.00 & \textbf{0.50} \\
\bottomrule
\end{tabular}
\end{table}

\textbf{Leave-one-out fails on real data.} On the synthetic workload it scored
$1.00$ and was separable from the anchored value only on the efficiency property.
Here, wherever the fault is smaller than the benign weekday/weekend shift, it
drops to $0.00$ alongside plain Shapley and stagewise, and the anchored value
alone holds at $1.00$ (Fig.~\ref{fig:main}b). Real drift is large and asymmetric
in a way our generator's was not: Saturday differs from Monday in trip distance,
hour distribution and tipping behaviour simultaneously, so removing a source from
the full coalition---exactly what leave-one-out does---moves the utility more
than the injected fault does. Anchoring removes that term by construction instead
of competing with it.

\textbf{Exclusions.} Only 6--8 of 19 configurations are valid test cases. Two
rules apply, both stated before results were read. First, \emph{the injected
fault causes no damage in isolation} (11 configurations): all four
\texttt{broken\_join\_key} levels and all three \texttt{join\_fanout} levels.
Duplicating rows in the zone lookup is textbook join fan-out, but a left join on
\texttt{DOLocationID} reproduces the same borough value, so it reweights the
training set uniformly and costs nothing; scoring a method for failing to blame
an operator that caused no harm would penalise it for being right. Second,
\emph{the fault is harmful but benign drift offset it} (2 configurations). The
second rule is itself a finding: on real data a genuine fault can be masked
entirely by a favourable partition, and a monitor watching only end-to-end model
quality would not fire at all.

\subsection{Scale}

A generated fan-in DAG with $m$ source branches ($n = 3m+3$) gives the cost
figures in Table~\ref{tab:scale}.

\begin{table}[t]
\caption{Cost as DAG width grows}
\label{tab:scale}
\centering
\begin{tabular}{rrrrrr}
\toprule
$n$ & $2^n$ & Fits & Cache hit & $|F|$ & MC time \\
\midrule
6  & 64            & 4  & 91.3\% & 2  & $<$0.01 s \\
9  & 512           & 8  & 98.5\% & 3  & $<$0.01 s \\
12 & 4{,}096       & 16 & 99.7\% & 4  & $<$0.01 s \\
15 & 32{,}768      & 3  & 99.4\% & 5  & 0.19 s \\
21 & $2.1\times10^6$ & 3 & 99.7\% & 7  & 0.41 s \\
27 & $1.3\times10^8$ & 3 & 99.8\% & 9  & 0.68 s \\
33 & $8.6\times10^9$ & 3 & 99.8\% & 11 & 0.99 s \\
\bottomrule
\end{tabular}
\end{table}

\begin{figure}[t]
\centering
\includegraphics[width=\columnwidth]{cost_scaling.pdf}
\caption{(a) The coalition lattice grows as $2^n$ while the number of model fits
stays flat. The step at $n=15$ is an artifact of the method, not the workload:
exact Shapley is computed up to $n=14$ and evaluates more coalitions than the
sampled estimator used beyond it. (b) The active frontier grows linearly and far
below $|V|$, which is what the flat fit count actually reflects.}
\label{fig:cost}
\end{figure}

\textbf{How to read this.} The headline ratio is real but is \emph{not} a
function of $n$ (Fig.~\ref{fig:cost}). Only three distinct sink states exist in
this workload, because every other operator is state-insensitive and its flip is
caught by the output hash before any model is fit. Cost tracks $|F|$, not $|V|$.
That identifies its own failure mode: the reduction collapses when a change
touches many operators simultaneously---a shared library upgrade, a global config
change, a backfill.

\section{Related Work}

\subsection{Debugging and Orchestrating ML Pipelines}

\textsc{mlinspect}~\cite{mlinspect,mlinspectj} extracts a dataflow DAG from a
pipeline and propagates lineage annotations to detect data-distribution bugs. It
is the closest prior art and is complementary: it inspects one execution, whereas
we compare two and attribute a quantity. Its DAG extraction is what a production
implementation of our method should reuse. \textsc{Modyn}~\cite{modyn} treats the
pipeline itself as the managed object, orchestrating data-centric retraining over
evolving datasets; it decides \emph{when} to retrain, where we ask why the last
retrain got worse. Recent position work on next-generation data engineering
pipelines~\cite{nextgen} argues that pipelines need first-class observability
semantics rather than bolt-on monitoring, which is the gap this paper addresses
for one specific question. Lineage systems~\cite{openlineage} supply the
topology our replay engine consumes but say nothing about magnitude.

\subsection{Shapley Values in Data Management}

\textsc{ShapleyPipe}~\cite{shapleypipe} applies Shapley values to pipeline
operators for \emph{construction}---which operator to select to maximise offline
accuracy. That is a design-time AutoML question over a space of candidate
pipelines; ours is a run-time forensic question about one deployed pipeline whose
inputs changed. The player sets, games and answers differ. A recent
survey~\cite{shapleysurvey} maps the broader use of Shapley values across data
analytics, and asymptotic analysis of the dataset-valuation
game~\cite{asymptotic} characterises its large-sample behaviour.

Efficient estimation is the shared bottleneck. \textsc{ShaplEIG}~\cite{shapleig}
selects which coalitions to evaluate by Bayesian experimental design, which is
directly applicable to our setting because our value function is expensive for
exactly the reason theirs is---each evaluation trains a model.
\textsc{OperatorSHAP}~\cite{operatorshap} accelerates estimation in a different
operator-valued setting. Our output-hash pruning is complementary to both: it
removes evaluations rather than choosing among them, and it composes with any
sampling scheme.

\subsection{Data Valuation}

Data valuation methods value \emph{training examples}; we value
\emph{operators}. The estimation machinery transfers, the semantics do not.
Recent work values data for LLM alignment via
Shapley~\cite{llmvaluation} and argues for influence-based auditing as a cheaper
alternative~\cite{beyondshapley}. The same efficiency-versus-fidelity tension
appears here, and is why we report exact values wherever brute force is
affordable.

\subsection{Data Quality and Agent Benchmarks}

Data-quality systems~\cite{deequ} verify declarative constraints, and
\textsc{SPADE}~\cite{spade} synthesises such assertions automatically for LLM
pipelines. Our benign controls (Table~\ref{tab:benign}) show that constraint
violation and model damage are distinct properties: a checker that never looks at
the model cannot separate them. Benchmarks such as
ELT-Bench~\cite{eltbench} and \textsc{DataGovBench}~\cite{datagovbench} evaluate
whether agents can \emph{author} or \emph{govern} pipelines; we ask which
operator in an existing pipeline broke.

\subsection{Root Cause Analysis}

Root-cause analysis for production systems has matured towards causal
formulations~\cite{causalrca} and agentic, evidence-grounded
reasoning~\cite{agenticrca}. Both target service-level incidents where the signal
is telemetry; our object is a dataflow DAG and our signal is model utility, which
admits the exact decomposition telemetry-based methods cannot offer. A further
body of recent work addresses ``AI-powered data pipeline observability,''
predominantly in preprint repositories and low-tier venues, without formal
guarantees, ground-truth benchmarks, or reproducible artifacts.

\subsection{Cooperative Game Theory}

The anchored value's restriction to feasible coalitions relates to conjunctive
permission structures~\cite{permission} and, when coalitions must be
DAG-connected, the Myerson value~\cite{myerson}. We state the connection without
proving a characterisation; doing so is the first open problem in
Sec.~\ref{sec:conclusion}.

\section{Limitations}\label{sec:limits}

\textbf{One real dataset, one pipeline shape.} Whether the crossover sits at the
same ratio for other domains, and for deep rather than wide pipelines, is
untested.

\textbf{The fault suite is weaker than intended.} Eleven of nineteen taxi
configurations were excluded because the injected fault does no damage, leaving
the real-data conclusion resting on eight test cases.

\textbf{Few valid cases at the crossover.} The regime separating the anchored
value from every baseline contains four configurations below $r=1$. The
separation is clean---$1.00$ against $0.00$---but it is four points, not a curve.

\textbf{The probe set must exist.} The method rests on a fixed evaluation set
representing ground truth. Where labels arrive with long delay, $u$ is not
computable as defined. This is the most serious practical limitation.

\textbf{Deterministic utility.} We use a deterministic learner with common random
numbers. For stochastic learners $u$ is a random variable, efficiency holds only
in expectation, and existing Shapley concentration bounds---which assume a
deterministic oracle---do not directly apply.

\textbf{Anchor choice is a modelling decision.} If a source change is itself the
repairable fault---a misconfigured extraction job rather than genuine world
drift---anchoring hides the culprit by construction. Detecting that case is
unsolved.

\textbf{Hashing requires determinism.} Operators involving sampling, timestamps,
or \texttt{ORDER BY} without a total order break output-stability pruning.

\textbf{A methodological note.} Three of our injected faults measured no damage
for reasons unrelated to the method: a pure rescale that downstream
standardisation undid, a label recomputed from the corrupted columns so the
target moved with the features, and a join fault against a feature the workload
gave no signal to. All three looked like plausible faults and none was. Any
fault-injection benchmark in this space must \emph{verify} that each injected
fault degrades the model before scoring attribution against it.

\section{Conclusion}\label{sec:conclusion}

Attributing model degradation to pipeline operators can be posed as a cooperative
game over hybrid replays of a fixed pipeline across two temporal snapshots. The
Shapley value decomposes degradation exactly but ranks poorly under co-occurring
benign drift. The incident-anchored Shapley value recovers both standard Shapley
and leave-one-out as boundary cases, ranks perfectly across every severity regime
tested, preserves exact decomposition, and---unlike leave-one-out---survives real
day-over-day drift.

Open problems, in the order we would attack them: an axiomatic characterisation
of $\phi^B$ and its precise relation to the conjunctive permission value;
automatic anchor selection, distinguishing genuine exogenous drift from a
repairable extraction fault; concentration bounds under a stochastic utility
oracle; the broad-change regime where the active frontier is large and pruning
fails; and federated attribution across data silos without centralising raw data.

\section*{Reproducibility}

All numbers are produced by the accompanying artifact; every experiment runs in
under two minutes on a laptop. The boundary theorems are asserted to machine
precision before any result is reported.

\begin{thebibliography}{00}

\bibitem{mlinspect} S. Grafberger, S. Schelter \emph{et al.}, ``mlinspect: a data distribution debugger for machine learning pipelines,'' in \emph{Proc. ACM SIGMOD}, 2021.

\bibitem{mlinspectj} S. Grafberger \emph{et al.}, ``Data distribution debugging in machine learning pipelines,'' \emph{The VLDB Journal}, vol. 31, 2022.

\bibitem{modyn} M. B\"other \emph{et al.}, ``Modyn: data-centric machine learning pipeline orchestration,'' arXiv:2312.06254, Dec. 2023.

\bibitem{nextgen} K. M. Kramer \emph{et al.}, ``Towards next generation data engineering pipelines,'' arXiv:2507.13892, Jul. 2025.

\bibitem{openlineage} OpenLineage, ``An open framework for data lineage collection and analysis.'' [Online]. Available: \url{https://openlineage.io}

\bibitem{shapleypipe} J. Chang, C. Liu, J. Huang, S. Zheng, R. Mao, and J. Qin, ``ShapleyPipe: hierarchical Shapley search for data preparation pipeline construction,'' arXiv:2510.27168, Oct. 2025.

\bibitem{shapleysurvey} H. Lin \emph{et al.}, ``A comprehensive study of Shapley value in data analytics,'' arXiv:2412.01460, Dec. 2024.

\bibitem{asymptotic} M. Tamine \emph{et al.}, ``An asymptotic analysis of the Shapley value for dataset valuation,'' arXiv:2607.03374, Jul. 2026.

\bibitem{shapleig} D. Rundel \emph{et al.}, ``ShaplEIG: Bayesian experimental design for Shapley value estimation,'' arXiv:2606.02247, Jun. 2026.

\bibitem{operatorshap} J. Stiller \emph{et al.}, ``OperatorSHAP: fast and accurate Shapley value estimation for neural operators,'' arXiv:2606.28065, Jun. 2026.

\bibitem{llmvaluation} M. Tamine \emph{et al.}, ``Shapley-based data valuation for LLM alignment via sequential preference optimization,'' arXiv:2512.15765, Dec. 2025.

\bibitem{beyondshapley} Y. Song \emph{et al.}, ``Beyond Shapley: an influence-based data auditing pipeline for LLM alignment and evaluation,'' arXiv:2607.22766, Jul. 2026.

\bibitem{deequ} S. Schelter \emph{et al.}, ``Automating large-scale data quality verification,'' \emph{Proc. VLDB Endowment}, vol. 11, no. 12, 2018.

\bibitem{spade} S. Shankar \emph{et al.}, ``SPADE: synthesizing data quality assertions for large language model pipelines,'' arXiv:2401.03038, Jan. 2024.

\bibitem{eltbench} ``ELT-Bench: an end-to-end benchmark for evaluating AI agents on ELT pipelines,'' \emph{Proc. VLDB Endowment}. doi:10.14778/3773749.3773750.

\bibitem{datagovbench} Z. Liu \emph{et al.}, ``DataGovBench: benchmarking LLM agents for real-world data governance workflows,'' arXiv:2512.04416, Dec. 2025.

\bibitem{causalrca} S. Jha \emph{et al.}, ``Causal AI-based root cause identification: research to practice at scale,'' arXiv:2502.18240, Feb. 2025.

\bibitem{agenticrca} A. Wei \emph{et al.}, ``Agentic root cause analysis through evidence-grounded reasoning,'' arXiv:2607.22385, Jul. 2026.

\bibitem{permission} R. van den Brink and R. P. Gilles, ``Axiomatizations of the conjunctive permission value for games with permission structures,'' \emph{Games and Economic Behavior}, vol. 12, no. 1, pp. 113--126, 1996.

\bibitem{myerson} R. B. Myerson, ``Graphs and cooperation in games,'' \emph{Mathematics of Operations Research}, vol. 2, no. 3, pp. 225--229, 1977.

\end{thebibliography}

\end{document}
